% ID: SEL-D0042-OB-E01
% ESTADO: validada. S043, cotejada con examen y criterios.
\textbf{a)} Sean $x,y$ los kg semanales de lácteos y pescado.
Minimizamos $C(x,y)=6x+12y$ euros con las restricciones
\[
x,y\ge0,\qquad x+y\le4,\qquad 3x+y\ge4,\qquad x+2y\ge3.
\]

\textbf{b)} La región queda por debajo de $y=4-x$ y por encima
de $y=0$, $y=4-3x$ e $y=(3-x)/2$, con $x\ge0$.
Los cortes factibles con los ejes son $(0,4)$, $(3,0)$ y $(4,0)$.
Las dos fronteras nutricionales se cruzan en $(1,1)$.
\begin{center}
\setlength{\unitlength}{9mm}
\begin{picture}(8,6)(-2,-1)
\qbezier(0,0)(2.35,0)(4.7,0)\put(4.8,0){$x$}
\qbezier(0,0)(0,2.2)(0,4.4)\put(0,4.5){$y$}
\qbezier(0,4)(2,2)(4,0)
\qbezier(0,4)(0.5,2.5)(1,1)
\qbezier(1,1)(2,0.5)(3,0)
\put(-0.2,4){\makebox(0,0)[r]{$(0,4)$}}
\put(1.1,1.1){$(1,1)$}
\put(2.8,-0.25){\makebox(0,0)[rt]{$(3,0)$}}
\put(4.1,-0.25){\makebox(0,0)[lt]{$(4,0)$}}
\end{picture}
\end{center}
Como $C=6(x+2y)$ y $x+2y\ge3$, el coste mínimo es
$\boxed{18}$ euros. Se alcanza en todo el lado
$y=(3-x)/2$ para $1\le x\le3$.
Por ejemplo, un kg de cada alimento y tres kg de lácteos sin
pescado son dos combinaciones óptimas distintas.
